\documentclass[10pt,a4paper]{amsart}
\usepackage[margin=19mm]{geometry}
\usepackage{amsmath,amssymb,mathtools}
\usepackage[scr=rsfs,cal=euler]{mathalfa}
\usepackage{xfrac}
\usepackage[unicode,hidelinks,bookmarks=false]{hyperref}
\newcommand{\ie}{i.e.\ }
\newcommand{\cf}{cf.\ }
\newcommand{\resp}{respectively\ }
\newcommand{\Subsection}{\subsection}
\providecommand{\nobreakdash}{\nobreak\hbox{-}}
\setlength{\emergencystretch}{2em}
\multlinegap0pt
\theoremstyle{plain}
\newtheorem{theorem}{Theorem}
\newtheorem{proposition}{Proposition}
\newtheorem{corollary}{Corollary}
\newtheorem{lemma}{Lemma}
\theoremstyle{definition}
\newtheorem{definition}{Definition}
\newtheorem{construction}{Construction}
\newtheorem{example}{Example}
\theoremstyle{remark}
\newtheorem{remark}{Remark}
\title[Hyperbolic links associated to Hamiltonian subgraphs in simple 3-polytopes]{Hyperbolic links associated to Hamiltonian subgraphs in simple 3-polytopes: links of Hamiltonian theta-subgraphs (printed as Theorem 6.3)}
\author{Nikolai Erokhovets}
\date{Source: arXiv:2512.03017v7 (CC BY 4.0) (CC BY 4.0)}
\begin{document}
\maketitle

\noindent\emph{Source and licence.} Hyperbolic links associated to Hamiltonian subgraphs in simple 3-polytopes. Version arXiv:2512.03017v7 (CC BY 4.0), distributed under the Creative Commons Attribution 4.0 International licence, http://creativecommons.org/licenses/by/4.0/, as recorded in the arXiv licence metadata for this exact version and shown on the saved abstract page https://arxiv.org/abs/2512.03017v7. Full licence notice: licenses/arxiv-2512.03017-CC-BY-4.0.txt.\par\smallskip

\noindent\textit{Editorial scope.} Counted unit: the classification of the link of a Hamiltonian theta-subgraph in a simple 3-polytope, printed in the article as Theorem 6.3 (source0.tex lines 1653--1663, one \texttt{theorem} environment with source label \texttt{Th:unl}), delivered together with the article's complete original proof of it (source0.tex lines 1664--1694, one \texttt{proof} environment of 31 lines immediately adjacent to the statement, with no gap and no deferral). The proof is the article's own geometric argument: it reads the linking of the circles off the representation of the theta-subgraph on the coordinate rays of the octant, produces a contradicting pair of matchings from any edge that joins two vertices of the same path, then shows that if the link is unlinked and the polytope is not the simplex then the pair is obtained by cutting off a vertex, and that otherwise three successive matching edges exhibit a triple of Borromean rings, which together with the article's Example 6.2 gives all three items of the classification. No mathematics is written, repaired or re-derived: statement and proof are the article's own text line for line, in the article's own notation, and no line of either is omitted, substituted or re-flowed.

Required context retained uncounted: the article's own Example 6.2 with its label (source0.tex lines 1646--1652), which records the trivial link of the theta-subgraph in the simplex and which the counted proof invokes in its last sentence. Nothing else from the article is needed for the statement or for the proof.

References. No citation command occurs in any retained line, so this package carries no bibliography block and no bibliography entry.

Editorial formatting only, in addition: an AMS \texttt{amsart} document, the same class the article itself uses, that restates the article's own theorem-environment declarations. No macro definition is needed by any retained line, so none is restated. Consequently the retained environments print here with independent counters, so the retained example prints as Example 1 and the counted theorem as Theorem 1, whereas the article prints them as Example 6.2 and Theorem 6.3 under its shared per-section counter; the equations of the retained spans are numbered anew in order of appearance and every cross-reference inside the retained text resolves to this wrapper's numbering, which is the only change to the numbering. No figure, no image-inclusion line and no drawing code occur in any retained line, so no artwork is delivered or needed, although the article contains figures elsewhere.

Not retained: the article's abstract, its constructions, its other theorems and examples, its figures and its bibliography; none of that material is used as a step of the delivered proof except the example that is retained, and the delivered proof is complete as the author wrote it. The package ships the article's concatenated source verbatim as \texttt{sources/190234/source0.tex}, and every delivered excerpt is an exact line range of that file. No mathematical statement, hypothesis, estimate or conclusion has been rewritten, repaired or added, and printed slips, if any, are preserved literally.
\begin{example}\label{ex:D30}
Let $\Gamma_0$ be the Hamiltonian theta-subgraph in the simplex $\Delta^3$, obtained by deletion of any edge 
from the graph $G(\Delta^3)=K_4$. Up to combinatorial symmetries it is a unique theta-subgraph in $G(\Delta^3)$.
The link $C_{\Gamma_0}$ is a trivial circle. Then for any pair $(P,\Gamma)$ obtained
from  $(\Delta^3,\Gamma_0)$ by a sequence of operations of cutting off a vertex the corresponding link~$C_{\Gamma}$
is trivial.
\end{example}

\begin{theorem}\label{Th:unl}
Let $\Gamma$ be a Hamiltonian theta-subgraph in a simple $3$-polytope $P$. Then 
\begin{enumerate}
\item the link $C_\Gamma$ consists of mutually unlinked circles if and only if each edge of $M_{\Gamma}$
connects vertices of different paths of $\Gamma$;
\item if $C_\Gamma$ consists of mutually unlinked circles and is nontrivial, then it contains 
a triple of Borromean rings; 
\item the link $C_\Gamma$ is trivial if and only if  $(P,\Gamma)$ is obtained from $(\Delta^3,\Gamma_0)$,
by a sequence of operations of cutting off a vertex.
\end{enumerate}
\end{theorem}

\begin{proof}[Proof of Theorem~\ref{Th:unl}]
From the representation of the theta-subgraph on the coordinate rays of the octant 
it is clear that each two circles are unlinked if each edge of the matching connects two different paths of $\Gamma$.
On the other hand, if there is an edge $M_\Gamma$
connecting two vertices $v$ and $w$ on the same path, then take such an edge $E_1$ 
with the condition that between $v$ and $w$ there are no pairs of vertices connected by edges in $M_\Gamma$. 
There is a vertex of another edge $E_2$ lying on the same path between $v$ and $w$, 
for otherwise there is a bigonal face, which is a contradiction. The edge $E_2$ lies 
in another connected component of $\partial P\setminus\Gamma$. The other vertex of $E_2$ lies
either on the same path, or on another path. In both cases it is clear from the octant representation
that the circles are linked is the standard way (as in the Hopf link). This proves~(1).

Now let $C_\Gamma$ consist of mutually unlinked circles. By (1) each edge of $M_{\Gamma}$ connects vertices on different paths
of $\Gamma$. Let $\Gamma_1$, $\Gamma_2$ and $\Gamma_3$ be the paths of $\Gamma$ connecting the vertices $v$ and $w$.
If there is an edge $E\in M_\Gamma$ with vertices $v_i$ and $v_j$ on $\Gamma_i$ and $\Gamma_j$ such that there are no vertices 
on these paths between $v$ and $v_i$ and $v$ and $v_j$,
then $P$ has a triangle incident to $v$ and if $P\ne \Delta^3$, then $(P,\Gamma)$ is obtained from 
some pair $(P',\Gamma')$ by cutting off a vertex.  The graph $G(P')$ is obtained from $G(P)$ by deletion of $E$. 
The link $C_{\Gamma}$ is trivial if and only if $C_{\Gamma'}$ is trivial. If $P$ has no such edges $E$,
then consider a vertex $v_1$ on $\Gamma_1$ closest to $v$. If this vertex is $w$, then all the edges of $M_{\Gamma}$ connect 
vertices on $\Gamma_2$ and $\Gamma_3$ and are ``parallel'', in particular the first and the last edges are of the above type.
A contradiction. Thus, $v_1$ belongs to some edge $E_1\in M_{\Gamma}$ with the other vertex $v_2$ 
lying on the other path, say $\Gamma_2$. By our assumption, there is a vertex $v_3$ between $v$ and $v_2$. Let $v_3$
be the closest vertex to $v$.  Then $v_3\in E_2\in M_\Gamma$. The other vertex $v_4$ of $E_2$ belongs to $\Gamma_3$.
Again by our assumption there is a vertex $v_5$ between $v$ and $v_4$, $v_5\in E_3\in M_\Gamma$. Let $v_6$
be the other vertex of $E_3$. Then $v_6\in \Gamma_1$ and $v_1$ lies between $v$ and $v_6$. The edges $E_1$,
$E_2$ and $E_3$ correspond to Borromean rings in $C_\Gamma$. In particular, $C_\Gamma$ is a nontrivial link.
Thus, if $C_\Gamma$ does not contain Borromean rings, then $(P,\Gamma)$ is obtained from 
$(\Delta^3,\Gamma_0)$ by a sequence of operations of cutting off a vertex. In particular, $C_\Gamma$ is trivial. Together
with Example~\ref{ex:D30} this proves (2) and (3).
\end{proof}

\end{document}
